\subsection*{Градиентный спуск}

Чтобы понять его суть, представьте себе горы в тумане. Видимость плохая, и ваша задача --- спуститься вниз. Что делать? Нужно смотреть под ноги и двигаться туда, где становится чуть ниже. Так шаг за шагом можно добраться до подножия. Принцип градиентного спуска примерно такой же: мы не знаем заранее, где находится минимум функции потерь, но можем на каждом шаге определять направление, в котором она убывает, и двигаться туда.

Это безусловно всё круто, но теперь осталось понять "--- а как нам найти это самое направление? Тут на помощь приходит понятие \textit{градиента} функции.

Давайте рассмотрим самую обыкновенную производную "--- как интерпретируется её значение? Вообще, численная величина производной отражает, насколько сильно изменяется функция в конкретной точке. В контексте обычной $f(x)$ это прокатит, но уже на $f(x, y)$ "--- всё рухнет. 

Если мы возьмём частную производную по $x$ (т.е. $\frac {\partial f}{\partial x}$), то получим <<насколько сильно изменится функция по $Ox$>>. Отдельно теперь возьмём производную по $y$. Объединяя эти значения в вектор, получим вектор, который указывает в сторону наискорейшего возрастания \textit{всей функции в целом}.

Итак, давайте чуть чуть формализуем. Представим, что у нас есть многомерная функция 
\[
f(x_1, \dots, x_n)
\]
её градиент обозначается и вычисляется следующим образом:
\[
\nabla f = (\frac{\partial f} {\partial x_1}, \dots, \frac{\partial f} {\partial x_n})
\]

Через $-\nabla f$ обозначается \textit{антиградиент}, и уже этот вектор указывает в сторону наискорейшего убывания функции. Кажется, ребус начинает потихоньку складываться.

И вот оно сходство с задачей спуска с горы "--- на каждом шаге мы спускаемся вниз и, с помощью антиградиента узнаём, куда нужно шагнуть на следующем шаге, чтобы оказаться ниже.

\newpage
Итак, выведем алгоритм градиентного спуска:
\begin{enumerate}
    \item Инициализируем вектор весов $w_0$ случайным образом;
    \item $w_k = w_{k-1} - \eta_k \nabla Q(x, w)$, где $Q$ есть выбранная функция потерь.
\end{enumerate}

$\eta_k$ --- шаг обучения или \textbf{learning rate}, который определяет, насколько далеко мы делаем каждый шаг (какую часть вектора антиградиента мы будем использовать для перемещения). $\eta_k \in [0, 1]$.

Выбор $\eta$ критически важен. Если шаг слишком большой, мы можем перепрыгивать через желанный минимум. Если шаг слишком маленький, обучение будет идти очень медленно или застопорится. 

На практике часто используют адаптивные стратегии: сначала шаг может быть большим, чтобы быстрее пройти начальные области, а затем уменьшается для более аккуратного приближения к минимуму. Именно поэтому у шага обучения указан индекс $k$ --- в зависимости от шага, он может изменяться.
Пока что мы упростим и полагаем, что $\eta_k$ --- определённая неизменяющаяся величина.

Полученный алгоритм и есть градиентный спуск, или полное его название --- \texttt{full batch gradient descent}. \texttt{Full batch} означает, что для одного изменения весов вычисляется ошибка на всём датасете сразу. Как Вы возможно догадались, именно эти дорогостоящие изменения весов мы попытаемся исправить далее.

\newpage
\subsection* {Стохастический Градиентный Спуск}

Вообще, в статистике есть такой трюк, что среднее всей выборки мы можем оценивать на каком"=то 
подмножестве примеров из этой выборки. Например, увидев 100 человек с ростом 175, вы с некоторой (признаю, довольно сильной) погрешностью можете полагать, что на Земле средний рост человека - 175. Это крайне грубый пример, но он позволит нам двигаться  дальше.

Идея \textbf{стохастического градиентного спуска} (\texttt{Stochastic Gradient Descent, SGD}) заключается в том, что мы будем оценивать 
\texttt {full batch gradient descent} 
не по какому"=то подмножеству из выборки, а... по одному примеру.

Давайте вот на что обратим внимание: аргумент суммы в функции потерь является некоторым преобразованием $x$. Значит, мы явно можем записать аргумент суммы в виде функции:
$$
Q(x, w) = \frac{1}{\ell}\sum_{i=1}^\ell q(x)
$$

Градиент --- линеен. Переводя на русский --- мы можем внести его в сумму:
$$
\nabla Q(x, w) = \frac{1}{\ell}\sum_{i=1}^\ell \nabla q(x)
$$

А теперь, чтобы значительно сократить количество вычислений, мы хотим
добиться следующего результата (то оценивание, о котором говорилось выше):
$$
\nabla Q(x, w) \approx \nabla q(x)
$$

И соответственно, формула принимает следующий облик:
\[
w_{k} = w_{k-1} - \eta \nabla_w q_{i_k}(w_{k-1})
\]
где $i_k$ "--- индекс случайного примера. 

\subsection*{Выбор learning rate}

Существует т.н. \textbf{условие Роббинса"=Монро}, которое гласит, что если \texttt{learning rate} удовлетворяет условиям:
\[
\sum_{i=1}^\infty \eta_i = \infty
\]
\[
\sum_{i=1}^\infty \eta_i^2 = const
\]
то градиентный спуск сойдётся с вероятностью 1.

Теперь давайте по"=человечески. Первая сумма ряда гарантирует, что спуск успевает дайти до минимума, т.е. шаг не замедляется слишком быстро. Вторая сумма ряда вносит некоторый баланс, и гарантирует, что при этом мы не можем уменьшать шаг слишком медленно.

\newpage
